% ====================================================================
% ================== Resumen en Español ==============================
% ====================================================================
\chapter{Resumen en Español}\label{chap:ResumenEspanol}

\section{Introducción}

Los vehículos autónomos avanzan hacia un despliegue generalizado, con el
potencial de mejorar de forma significativa la seguridad vial, aliviar la
congestión y aumentar la eficiencia global del transporte. En el núcleo de esta
transformación se encuentra la capacidad de percibir el entorno con elevada
precisión y robustez. Los sistemas de percepción permiten a los vehículos
detectar obstáculos, interpretar el comportamiento de otros agentes del tráfico
y comprender el contexto general de la escena de conducción. Sin embargo, el
mundo real es dinámico y a menudo poco estructurado, lo que plantea importantes
desafíos a los algoritmos de percepción.

Los fallos de percepción pueden desembocar en \textit{disengagements}, es
decir, situaciones en las que el sistema automatizado ya no puede identificar
un curso de acción seguro y debe transferir el control a un operador humano.
Estos eventos se producen habitualmente fuera de las condiciones óptimas de
funcionamiento del sistema y resultan especialmente críticos. Incluso en tales
escenarios, el vehículo debe mantener un cierto nivel de conocimiento del
entorno, de forma que se posibilite la activación de estrategias de
\textit{fallback} o se asista al operador remoto para resolver la situación de
manera segura.

Los sistemas de percepción más avanzados en la actualidad suelen basarse en
técnicas de aprendizaje profundo, entrenadas con grandes volúmenes de datos
anotados. Aunque estos métodos alcanzan un alto rendimiento en entornos bien
estructurados y representados en los conjuntos de entrenamiento, tienden a
mostrar problemas de generalización en situaciones raras o previamente no
observadas. Su falta de transparencia plantea además interrogantes en torno a
la interpretabilidad y la validación sistemática, aspectos clave en
aplicaciones de seguridad crítica.

Para afrontar estas limitaciones, esta tesis explora un enfoque complementario
y determinista basado en información de cámaras estéreo. El objetivo es diseñar
e implementar un sistema de percepción capaz de estimar la posición y el
movimiento de los objetos circundantes sin depender de entrenamiento previo ni
de anotaciones semánticas. Al centrarse en pistas geométricas, como la
disparidad y el flujo óptico, el sistema busca un rendimiento robusto en
entornos diversos y desafiantes.

El sistema procesa imágenes estéreo sincronizadas captadas por un conjunto de
cámaras embarcadas para producir estimaciones de movimiento en seis
dimensiones: tres posiciones espaciales y las correspondientes tres
componentes de velocidad de puntos seleccionados de la escena. Un marco basado
en filtros de Kalman, inspirado en el enfoque \textbf{6D Vision} de Daimler
\cite{6d_vision_daimler}, integra profundidad, flujo óptico y egomoción para
calcular dichas estimaciones. El campo de movimiento resultante se agrupa
mediante el algoritmo \textbf{DBSCAN}, que permite identificar objetos móviles
coherentes.

Con el fin de evaluar posibles riesgos, se desarrolla además un módulo de
criticidad que calcula métricas de seguridad a partir de las trayectorias
predichas tanto del ego-vehículo como de los objetos detectados. Una de las
métricas más relevantes es el \textbf{Tiempo hasta Colisión} (\textit{Time-to-
Collision}, TTC), que ofrece información directa sobre la urgencia de
maniobras evasivas. Estas evaluaciones se realizan en tiempo real, de manera
que el sistema pueda servir de apoyo tanto a estrategias de \textit{fallback}
autónomo como a la toma de decisiones en teleoperación.

El trabajo se apoya en investigaciones previas realizadas en la Cátedra de
Tecnología del Automóvil de la Universidad Técnica de Múnich, centradas en
visión estéreo y estimación determinista de movimiento. La presente tesis se
concentra en integrar estos métodos en un componente software modular y
reutilizable. Una contribución destacada es el desarrollo de un servicio
independiente en \textbf{ROS2}, diseñado para su integración fluida en el
vehículo de investigación EDGAR y adaptable a otras plataformas.

El sistema se implementa en C++, con componentes críticos en cuanto a
rendimiento acelerados mediante \textbf{CUDA} y paralelizados con \textbf{OpenMP}.
De esta manera, se cumplen las restricciones de tiempo real al mismo tiempo que
se asegura escalabilidad y robustez. Cada función principal, ya sea la
estimación de profundidad, el cálculo de flujo óptico, el filtrado de Kalman,
el agrupamiento o la evaluación de criticidad, se encapsula en un paquete ROS2
independiente, lo que facilita el desarrollo modular, las pruebas y la
integración. El uso de Docker para contenerización aporta portabilidad y
simplifica el despliegue en sistemas heterogéneos.

La arquitectura del sistema se compone de dos pipelines principales:

\begin{itemize}
  \item \textbf{Pipeline de percepción:} combina visión estéreo, odometría
  visual y estimación determinista de movimiento para inferir trayectorias en
  seis dimensiones de puntos de la escena. Los objetos en movimiento se
  identifican agrupando estas trayectorias en función de su coherencia espacial
  y dinámica.

  \item \textbf{Pipeline de criticidad:} analiza la interacción entre el ego-
  vehículo y los objetos detectados para calcular métricas de riesgo que
  apoyen decisiones de \textit{fallback} o asistencia al conductor.
\end{itemize}

La evaluación del sistema se organiza en tres fases. Primero, se calibra y se
valida el hardware estéreo para asegurar su idoneidad en estimaciones 6D de
movimiento; si es necesario, se integra hardware alternativo. En segundo lugar,
se evalúa el rendimiento computacional de todo el pipeline de percepción y
criticidad para verificar su capacidad de operar en tiempo real. Por último, se
despliega el sistema en el vehículo EDGAR para su validación en condiciones
reales, donde sus salidas se comparan con las de la pila de percepción basada
en aprendizaje de Autoware. Todos los experimentos se diseñan para comprobar el
cumplimiento de los requisitos funcionales definidos en el
Capítulo~\ref{sec:RequirementsForTheThesis}, que incluyen la restricción de
usar únicamente estéreo, la independencia semántica, la modularidad, la
operación en tiempo real y la robustez frente a oclusiones.

En síntesis, la tesis presenta un sistema determinista, modular y capaz de
operar en tiempo real para la percepción en conducción automatizada. Al
aprovechar la visión estéreo y algoritmos geométricos, el sistema amplía la
capacidad de comprensión del entorno en situaciones donde los métodos de
aprendizaje profundo pueden fallar. Su despliegue como servicio en ROS2,
optimizado para rendimiento y portabilidad, permite una integración sencilla en
entornos de investigación y desarrollo diversos.

Finalmente, la estructura del documento es la siguiente: el
Capítulo~\ref{chap:Introduction} introduce la motivación, los objetivos y el 
planteamiento general del trabajo. El Capítulo~\ref{chap:StateOfTheArt} revisa 
los sistemas de percepción actuales para conducción automatizada, incluyendo 
sensores como LiDAR, RADAR y cámaras. Se analizan también algoritmos para 
detección de objetos móviles, diferenciando entre enfoques de aprendizaje y 
deterministas, y se resumen los trabajos relevantes realizados en la Cátedra de 
Tecnología del Automóvil. El Capítulo~\ref{chap:MethodsAndImplementation} 
describe el diseño e implementación del sistema propuesto, separando las 
pipelines de percepción y criticidad en módulos diferenciados. El 
Capítulo~\ref{chap:ExperimentsAndResults} presenta la evaluación, tanto a nivel 
de componentes como en pruebas reales con el vehículo EDGAR. Finalmente, el 
Capítulo~\ref{chap:SummaryAndConclusion} resume las principales contribuciones 
y plantea futuras líneas de trabajo.


\section{Estado del arte}

La investigación en percepción para vehículos automatizados se ha centrado en
integrar diferentes modalidades sensoriales para obtener una representación
fiable del entorno. Los sistemas actuales combinan habitualmente sensores como
LiDAR, radar y cámaras, aprovechando las ventajas particulares de cada uno. El
LiDAR ofrece mapas tridimensionales de alta precisión, aunque es costoso y
sensible a condiciones meteorológicas adversas. El radar resulta robusto bajo
lluvia o niebla, y permite medir velocidades relativas gracias al efecto
Doppler, si bien su resolución espacial es limitada. Las cámaras, por su parte,
aportan información visual rica, de bajo coste y con gran densidad, aunque
dependen fuertemente de la iluminación y presentan dificultades en regiones con
poca textura.

Además de la elección del sensor, los algoritmos de detección de objetos en
movimiento pueden clasificarse en dos grandes familias: los enfoques basados en
aprendizaje y los enfoques deterministas. Los primeros, apoyados en redes
neuronales convolucionales o arquitecturas más recientes como transformadores,
han alcanzado resultados destacados en competiciones y bases de datos públicas.
Su limitación principal reside en la necesidad de grandes volúmenes de datos
anotados y en la dificultad para generalizar a situaciones poco representadas.
Los segundos, de carácter determinista, se fundamentan en principios físicos y
geométricos. Entre ellos destacan el cálculo de flujo óptico, la odometría
visual y el uso combinado de disparidad y movimiento para inferir trayectorias.
El enfoque denominado \textbf{6D Vision} \cite{6d_vision_daimler} pertenece a
esta última categoría y plantea estimar, para cada punto, un vector de estado
con posición y velocidad. Esta técnica se ha mostrado útil en contextos donde
la semántica es secundaria y lo prioritario es disponer de medidas geométricas
robustas y explicables.

En la Cátedra de Tecnología del Automóvil de la TUM se han desarrollado varios
trabajos previos relacionados con este paradigma, explorando su viabilidad en
teleoperación y en conducción automatizada. Sin embargo, hasta ahora las
implementaciones habían mostrado limitaciones de rendimiento en tiempo real o
dificultades de integración con la arquitectura de software del vehículo
EDGAR. Esta tesis recoge estas experiencias previas y propone una
implementación modular, acelerada en GPU y desplegable en ROS2, con el objetivo
de superar las limitaciones detectadas y evaluar la utilidad del enfoque
determinista en comparación directa con un sistema de referencia basado en
aprendizaje como Autoware.


\section{Pipeline de percepción}

La \textbf{pipeline de percepción} constituye el núcleo del sistema. Su
objetivo es transformar imágenes estéreo sin procesar en una representación
dinámica del entorno en la que cada objeto circundante se describa por su
posición tridimensional y su vector de velocidad. Para ello, se emplean varias
etapas que trabajan de forma encadenada: cálculo de flujo óptico, estimación de
odometría visual, reconstrucción de profundidad, integración de información en
un filtro de Kalman en seis dimensiones y, finalmente, agrupamiento de puntos
en entidades coherentes mediante clustering.

El flujo óptico proporciona información sobre el desplazamiento aparente de los
píxeles entre imágenes consecutivas. Su formulación clásica parte de la
ecuación de constancia de brillo, que asume que la intensidad de un punto se
mantiene aproximadamente constante a lo largo del tiempo. De esta premisa se
deduce la conocida ecuación diferencial que relaciona las derivadas espaciales
de la intensidad con las componentes del flujo. A lo largo de la literatura se
han propuesto distintos métodos de resolución, desde los de tipo local, como el
de Lucas–Kanade \cite{opt_flw_lukas_canade}, hasta los de tipo global, como el
de Horn–Schunck \cite{opt_flw_schunck}. En este trabajo se emplea el algoritmo
de Farnebäck \cite{opf_flw_farneback}, que aproxima cada vecindad de la imagen
mediante un polinomio cuadrático y obtiene un campo denso de desplazamientos.
La implementación en GPU mediante CUDA permite reducir de manera notable los
tiempos de cómputo, alcanzando frecuencias superiores a 10 Hz, requisito clave
para la operación en tiempo real.

La segunda etapa del pipeline es la estimación de la egomoción mediante
odometría visual estéreo. Este módulo calcula la transformación relativa entre
pares de imágenes consecutivas, de modo que se distingue entre el movimiento
propio del vehículo y el movimiento de los objetos externos. El procedimiento
consiste en detectar puntos característicos (ORB), extraer descriptores
binarios, emparejarlos entre fotogramas y reconstituir su posición tridimensional
a partir de la disparidad del frame previo. Con estas correspondencias 3D–2D se
resuelve un problema de \textit{Perspective-n-Point} (PnP) con un esquema de
robustificación mediante RANSAC. El resultado es la rotación y traslación
estimada del sistema de cámaras. Para suavizar las oscilaciones entre
estimaciones sucesivas, se aplica un filtrado exponencial. Además, se estima de
forma aproximada la covarianza de la egomoción a partir de los residuales de
reproyección, lo que permite incorporar incertidumbre explícita en la etapa de
fusión posterior.

La tercera pieza es la reconstrucción de profundidad, que se basa en los
principios de la geometría estéreo. En un sistema rectificado, la disparidad
entre las posiciones de un mismo punto en la imagen izquierda y derecha es
inversamente proporcional a la profundidad. Matemáticamente se expresa como

\[
Z = \frac{f \cdot B}{d}
\]

donde $Z$ es la profundidad, $f$ la distancia focal, $B$ la línea base y $d$ la
disparidad. Para calcular este último valor se utiliza un algoritmo de tipo
Semi-Global Matching (SGM), que combina criterios de correspondencia local con
regularización global, logrando mapas de disparidad densos y suficientemente
robustos para escenas complejas. Estos mapas se proyectan al espacio tridimensional
y sirven de entrada al filtrado en seis dimensiones.

El núcleo del pipeline se inspira en el paradigma \textbf{6D Vision}
\cite{6d_vision_daimler}. Este concepto propone asociar a cada punto de la
escena un vector de estado compuesto por tres componentes de posición y tres de
velocidad. La estimación se realiza mediante un filtro de Kalman, que integra
información de disparidad, flujo óptico y egomoción. El modelo de transición se
basa en movimiento constante, expresado como

\[
\mathbf{x}_{k|k-1} = \mathbf{F}\mathbf{x}_{k-1|k-1} + \mathbf{w}_k,
\]

donde $\mathbf{F}$ es la matriz de transición que incorpora la relación entre
posiciones y velocidades y $\mathbf{w}_k$ representa el ruido de proceso. La
fase de observación utiliza las posiciones derivadas de la disparidad y las
velocidades proyectadas del flujo óptico, según

\[
\mathbf{z}_k = \mathbf{H}\mathbf{x}_{k|k-1} + \mathbf{v}_k,
\]

con $\mathbf{H}$ como matriz de observación y $\mathbf{v}_k$ el ruido de
medida. Adicionalmente, la estimación de egomoción obtenida en el bloque
anterior se emplea para compensar el movimiento propio del vehículo y expresar
las trayectorias en un marco de referencia estable. De este modo, el sistema es
capaz de separar el desplazamiento de la cámara de los movimientos reales de la
escena.

El resultado de este proceso es un conjunto de trayectorias 6D para miles de
puntos de la escena. No obstante, trabajar únicamente con puntos individuales
dificulta la interpretación a nivel de objetos. Por ello, se añade una etapa de
agrupamiento mediante el algoritmo DBSCAN \cite{dbscan_original}. Este método
tiene la ventaja de no requerir conocer a priori el número de objetos y de ser
robusto frente a ruido. Se define una métrica híbrida que combina distancia
espacial y coherencia de velocidad:

\[
d(a,b) = w_p \|p_a - p_b\| + w_v \theta(\vec{v}_a, \vec{v}_b),
\]

donde $\theta$ es el ángulo entre vectores de velocidad y $w_p$, $w_v$ son
pesos configurables. Así, dos puntos sólo se agrupan si están próximos en el
espacio y se mueven de forma coherente. Una vez formados los clusters, cada
conjunto se representa como un objeto mediante el cálculo de su centroide,
velocidad media y una caja delimitadora orientada. Esta última se obtiene a
través de un análisis de componentes principales (PCA), que permite alinear la
caja con la dirección predominante de los puntos. El resultado son entidades
\textit{WorldEntity} que constituyen la interfaz hacia el módulo de criticidad.

En conjunto, la pipeline de percepción transforma señales crudas en una
descripción de alto nivel del entorno. Parte de imágenes estéreo, estima
movimiento 6D a nivel puntual, filtra trayectorias con un modelo dinámico y
finalmente abstrae la información en objetos. El diseño modular y la
implementación en C++ con paralelización en GPU y CPU garantizan que todo este
proceso pueda realizarse en tiempo real, requisito indispensable para su
integración en vehículos automatizados.


\section{Pipeline de criticidad}

Una vez que la pipeline de percepción ha transformado las imágenes estéreo en
una lista estructurada de objetos con posición y velocidad, resulta necesario
evaluar el riesgo que cada uno de estos elementos supone para el ego-vehículo.
Este análisis constituye la función de la \textbf{pipeline de criticidad}, cuyo
objetivo es traducir la información geométrica en indicadores de seguridad que
puedan emplearse en la toma de decisiones, tanto en un modo de conducción
automatizada como en teleoperación.

El concepto central que se utiliza en esta tesis es el \textbf{Tiempo hasta
Colisión} (\textit{Time-to-Collision}, TTC). El TTC se define como el intervalo
temporal que transcurriría hasta que el vehículo y un objeto ocupasen la misma
posición si se mantienen sus trayectorias actuales. Se trata de una métrica
extendida en la literatura de conducción automatizada
\cite{cricitality_metrics}, ya que proporciona un indicador intuitivo y directo
del grado de urgencia de una maniobra evasiva. A diferencia de otras métricas
como la distancia temporal de seguridad o el índice de exposición al riesgo, el
TTC tiene la ventaja de poder calcularse de forma sencilla a partir de
información de posición y velocidad.

El cálculo del TTC puede formularse de distintas maneras. En una primera
aproximación, se considera la proyección de la velocidad relativa sobre la
línea de unión entre el ego-vehículo y el objeto. Si se denota por $d$ la
distancia actual y por $v_{rel}$ la componente de la velocidad relativa en esa
dirección, el TTC puede expresarse como

\[
TTC = \frac{d}{v_{rel}}.
\]

Este enfoque, denominado en esta tesis \textit{twist-based}, es adecuado cuando
tanto el vehículo como el objeto se desplazan en trayectorias rectilíneas y
constantes. Sin embargo, en escenarios urbanos o en maniobras complejas, esta
hipótesis resulta insuficiente. Por ello se implementa un segundo modo de
cálculo, denominado \textit{path-based}, que considera explícitamente la
trayectoria planificada del vehículo. En este caso, el TTC se define como el
primer instante en el que la trayectoria del ego-vehículo interseca con la del
objeto, teniendo en cuenta posiciones, velocidades y geometrías de colisión. En
su versión más simple, el cálculo se realiza sobre representaciones puntuales,
pero el sistema está diseñado para extenderse a colisiones entre cajas
delimitadoras orientadas.

La pipeline de criticidad recibe como entradas la lista de objetos detectados
por el módulo de percepción y la trayectoria prevista del ego-vehículo. Esta
última puede derivarse de la odometría o de un planificador de ruta. A partir
de estas entradas, se calculan en paralelo ambas variantes del TTC. Los
resultados se publican en tópicos ROS2, tanto a nivel de cada objeto como en
forma de un mínimo global, que representa el tiempo hasta la colisión más
inminente.

Un aspecto relevante del diseño es que la pipeline de criticidad se mantiene
independiente de la fuente de percepción. Para comparar el sistema propuesto
con Autoware, se implementó un adaptador que traduce las salidas de Autoware al
formato de objetos definido en esta tesis. De esta manera, ambos pipelines de
percepción pueden ser evaluados bajo las mismas métricas de criticidad, lo que
asegura una comparación justa y coherente. En la práctica, esto significa que
las discrepancias observadas en los resultados de TTC reflejan diferencias en
la percepción subyacente y no en la metodología de evaluación.

El diseño modular de este bloque ofrece flexibilidad adicional. Por ejemplo,
permite sustituir la métrica TTC por otras métricas de riesgo descritas en la
literatura, como el \textit{Post-Encroachment Time}, el \textit{Crash
Imminence Metric} o el cálculo de estados de colisión inevitables. Así, la
pipeline de criticidad no sólo cumple la función de evaluar la seguridad en
tiempo real, sino que también constituye una plataforma experimental para
explorar distintos indicadores de riesgo y estudiar su comportamiento en
condiciones reales.

En síntesis, la pipeline de criticidad convierte las estimaciones geométricas
del pipeline de percepción en indicadores cuantitativos de riesgo. Su
implementación en ROS2, independiente de la fuente de percepción, garantiza la
reutilización y la comparación con sistemas alternativos. El uso del TTC como
métrica central proporciona una herramienta clara para priorizar objetos y
determinar la urgencia de maniobras evasivas, tanto en contextos de conducción
autónoma como en situaciones de teleoperación.


\section{Evaluación experimental}

La validación del sistema desarrollado se llevó a cabo siguiendo una
metodología progresiva. Primero se abordó la selección y calibración del
hardware de captura. Después se realizaron pruebas controladas en escenarios
delimitados, donde era posible medir con precisión las trayectorias de los
objetos. Finalmente, se efectuaron ensayos en conducción urbana real con el
vehículo EDGAR, ejecutando el sistema determinista en paralelo con la pila de
percepción de Autoware. De esta forma se logró una evaluación tanto cualitativa
como cuantitativa de las prestaciones alcanzadas.

\bigskip
\noindent \textbf{Selección y calibración de cámaras.}  
Se evaluaron dos cámaras estéreo disponibles en el laboratorio: la Framos D455e
y la Carnegie S30. La Framos es una cámara de tipo activo, con proyección de
infrarrojos integrada, mientras que la Carnegie es de tipo pasivo con una base
estereoscópica de 30 cm. En ambos casos, se llevó a cabo una calibración
intrínseca mediante un patrón de tablero de ajedrez y se calcularon parámetros
de distorsión radial y tangencial. La calibración extrínseca se realizó
empleando técnicas de alineación con marcadores visibles, asegurando que las
matrices de rotación y traslación entre las dos cámaras alcanzasen errores
subpíxel. Además, se validó la compatibilidad con el protocolo de
\textit{Precision Time Protocol} (PTP) para garantizar la sincronización de las
imágenes con el resto de sensores de EDGAR.

\begin{table}[H]
  \centering
  \begin{tabular}{|l|c|c|}
    \hline
    \textbf{Cámara} & \textbf{Framos D455e} & \textbf{Carnegie S30} \\ \hline
    Tipo & Estéreo activo (IR) & Estéreo pasivo \\
    Baseline & 9 cm & 30 cm \\
    Alcance útil & $\sim$8 m & >15 m \\
    Robustez en exteriores & Limitada & Alta \\
    Ruido en baja textura & Medio/alto & Bajo \\
    Compatibilidad PTP & Parcial & Completa \\
    Observaciones & Artefactos en oblicuo & Mejor estabilidad en flujo \\ \hline
  \end{tabular}
  \caption{Comparación de las cámaras evaluadas para el sistema}
  \label{tab:camera_comparison}
\end{table}

Los resultados de esta primera fase mostraron que la Framos tenía un alcance
útil inferior y una mayor sensibilidad a artefactos en trayectorias oblicuas,
con aparición de sombras de disparidad. La Carnegie S30, en cambio, ofreció
mapas de profundidad más estables, coherentes con los resultados reportados en
\cite{paperCarnegie}, y fue seleccionada como cámara principal para los
experimentos.

\bigskip
\noindent \textbf{Pruebas controladas.}  
Para evaluar el rendimiento básico del sistema, se diseñaron escenarios
específicos. En uno de ellos, un peatón cruzaba lateralmente a una distancia
constante de ocho metros. El sistema detectó el objeto de manera continua,
generando un cluster estable y velocidades alineadas con la cinemática real. En
otro escenario, un objeto se desplazaba de forma oblicua con un ángulo de 18°
respecto al eje de visión. El sistema fue capaz de identificar la dirección de
velocidad sin desviaciones notables, demostrando la eficacia de la compensación
de egomoción. Estas pruebas se repitieron varias veces bajo diferentes
condiciones de iluminación, incluyendo entornos crepusculares, con el fin de
evaluar la robustez. En todos los casos, la Carnegie S30 mostró un nivel bajo
de ruido en regiones homogéneas, lo que facilitó la coherencia de las
estimaciones.

\bigskip
\noindent \textbf{Ensayos urbanos y comparación con Autoware.}  
La fase más exigente de la validación consistió en desplegar el sistema en el
vehículo EDGAR y recorrer entornos urbanos de media complejidad. Se ejecutaron
en paralelo tanto el pipeline determinista como la pila de Autoware, alojados en
contenedores Docker separados. Para garantizar una comparación justa, se
implementó un adaptador que tradujo los objetos de Autoware al formato de
\texttt{ClusteredObjectArray}, de forma que el cálculo de \textit{Time-to-
Collision} (TTC) se realizó con la misma metodología para ambos sistemas.

Los resultados mostraron diferencias relevantes. En un escenario de cruce de
peatones, el pipeline determinista identificó la trayectoria de un individuo
con ropa oscura en una zona de bajo contraste, mientras que Autoware no generó
una detección válida. En cambio, en un entorno con varios vehículos, Autoware
produjo una clasificación semántica más detallada, mientras que el sistema
determinista agrupó parte de la escena en un único cluster, perdiendo
granularidad. Este resultado confirma que ambos enfoques son complementarios:
uno aporta interpretabilidad y robustez, el otro riqueza semántica y
completitud.

\bigskip
\noindent \textbf{Rendimiento computacional.}  
Las pruebas de rendimiento se llevaron a cabo en la unidad GPU del EDGAR. El
sistema alcanzó frecuencias promedio de 12 Hz, superando el requisito mínimo de
10 Hz. El cálculo de flujo óptico en CUDA redujo la carga de CPU en un 35\% en
comparación con la versión puramente en CPU, mientras que la paralelización con
OpenMP en el filtro de Kalman permitió escalar el número de puntos procesados
sin degradar la frecuencia. En términos de recursos, la utilización promedio de
CPU fue del 45\% y la de GPU del 60\%, manteniéndose en márgenes compatibles
con el resto de procesos del vehículo. La latencia end-to-end desde la captura
de la imagen hasta la publicación del TTC fue inferior a 90 ms, valor adecuado
para integrar el sistema en un lazo de control vehicular.

\bigskip
\noindent \textbf{Limitaciones observadas.}  
Durante la evaluación se detectaron algunos aspectos susceptibles de mejora. En
primer lugar, el clustering tendía a fusionar objetos cercanos con velocidades
similares, lo que podía simplificar en exceso la representación de la escena.
En segundo lugar, el cálculo del TTC en modo rectilíneo mostraba inestabilidad
en trayectorias no lineales, lo que confirma la necesidad de extender el método
a colisiones basadas en cajas delimitadoras orientadas. Por último, la
resolución de la cámara Carnegie limitaba la precisión en objetos muy pequeños
o alejados, lo que abre la posibilidad de incorporar técnicas de submuestreo
adaptativo o fusión con otros sensores.

\bigskip
En síntesis, la evaluación experimental demostró que el sistema propuesto es
capaz de operar en condiciones reales, detectar y seguir objetos con coherencia
temporal y calcular métricas de criticidad en tiempo de ejecución. La
comparación con Autoware reveló la complementariedad de ambos enfoques y puso
de relieve la importancia de los métodos deterministas como capa de redundancia
y seguridad en conducción automatizada.



\section{Conclusiones y trabajo futuro}

El trabajo desarrollado en esta tesis ha tenido como objetivo el diseño e
implementación de un sistema determinista de percepción basado exclusivamente
en cámaras estéreo. A lo largo de las distintas fases se ha demostrado que es
posible alcanzar una operación en tiempo real, con modularidad suficiente para
sustituir componentes de forma independiente y con independencia semántica,
evitando la necesidad de entrenamiento previo. El sistema cumple, por tanto,
con los requisitos planteados al inicio: procesar en línea información visual,
distinguir entre movimiento propio del vehículo y movimiento externo, agrupar
trayectorias en objetos coherentes y calcular métricas de criticidad útiles para
la toma de decisiones en conducción automatizada o en teleoperación.

Las aportaciones técnicas pueden resumirse en varios ejes. En primer lugar, la
modernización del paradigma \textbf{6D Vision} mediante una implementación
modular en ROS2, con soporte de aceleración en GPU y CPU. En segundo lugar, la
incorporación de un algoritmo de clustering robusto como DBSCAN con métrica
híbrida, que permite agrupar puntos en objetos sin necesidad de conocer su
número a priori. En tercer lugar, la integración de una \textbf{pipeline de
criticidad} que calcula en tiempo real métricas como el \textit{Time-to-
Collision}, reutilizable tanto para el sistema desarrollado como para otros
pipelines, como Autoware. Finalmente, la validación práctica en el vehículo
EDGAR, que ha mostrado la viabilidad del despliegue en un entorno real con
restricciones de hardware y sincronización estrictas.

Pese a estos logros, el trabajo presenta algunas limitaciones que deben
considerarse. Entre ellas, destacan: 
\begin{itemize}
  \item El modelo de proceso del filtro de Kalman, que se basa en hipótesis de
  movimiento constante y no incorpora aceleraciones ni modelos dinámicos más
  complejos.
  \item La ausencia de un algoritmo avanzado de asociación multiobjeto, lo que
  puede limitar la persistencia de los objetos en escenas densas o en presencia
  de oclusiones prolongadas.
  \item El cálculo del TTC, que en su versión actual se limita al análisis de
  distancias entre centros de caja, en lugar de considerar colisiones entre
  cajas delimitadoras orientadas (OBB).
\end{itemize}

Estas limitaciones abren líneas claras de trabajo futuro. Una primera dirección
consiste en integrar información adicional procedente de sensores como IMU o
LiDAR, lo que aumentaría la robustez y ampliaría el rango operativo. Una
segunda dirección es la implementación de algoritmos de asociación de objetos
entre frames, por ejemplo mediante técnicas de \textit{tracking} multiobjetivo
con estrategias de \textit{data association} y modelos de predicción más
avanzados. Una tercera línea se centra en extender el cálculo del TTC para
considerar geometrías completas de colisión, como la intersección entre OBB,
lo que permitiría obtener estimaciones más realistas y precisas en entornos
urbanos complejos. Finalmente, se podrían explorar métricas adicionales de
criticidad, tales como el \textit{Post-Encroachment Time} o indicadores basados
en estados de colisión inevitables, con el fin de enriquecer la evaluación de
riesgo en tiempo real.

Más allá de estas propuestas, cabe destacar la relevancia del enfoque
determinista en un contexto más amplio. En un ecosistema dominado por métodos
de aprendizaje profundo, disponer de una alternativa transparente y explicable
es de gran valor para la seguridad funcional. Normativas como ISO~26262 o SOTIF
exigen argumentos claros de seguridad y verificación, algo que este sistema
puede proporcionar gracias a su naturaleza basada en principios geométricos y
modelos matemáticos bien definidos. Además, su carácter modular y portable
facilita la integración como capa de redundancia en arquitecturas de percepción
existentes.

En conjunto, el trabajo demuestra que un sistema determinista basado en cámaras
estéreo puede complementar de manera eficaz a las soluciones de percepción
aprendidas, aportando transparencia, robustez y explicabilidad en escenarios
críticos. Su implementación en un vehículo real confirma la viabilidad práctica
del enfoque y abre la puerta a futuras investigaciones orientadas a combinar lo
mejor de ambos mundos: la riqueza semántica del aprendizaje profundo y la
fiabilidad verificable de los métodos deterministas.
